\documentclass[a4paper, 11pt]{article}          % Paper and font size
\usepackage[english]{babel}                     % Language setting

% Set margins
\usepackage[top=2.5cm,bottom=2.5cm,left=2cm,right=2cm,marginparwidth=1.75cm]{geometry}

% Load useful packages
\usepackage{graphicx}                           % Graphics
\usepackage{amsmath}                            % Math
\usepackage{xcolor}                             % Link color
\definecolor{custom-blue}{RGB}{0,99,166}
\usepackage{hyperref}
\hypersetup{colorlinks=true, allcolors=custom-blue}
\usepackage[default]{sourcesanspro}             % Load font
\usepackage[T1]{fontenc}                        % Special characters
\usepackage{fancyhdr}                           % Load header, footer package
\usepackage{titlesec}

\usepackage{lastpage}                           % Footer note

\usepackage[table,xcdraw]{xcolor}               % Table requirements
\usepackage{listings}                           % Code listings
\usepackage{booktabs}                           % Nicer tables

\definecolor{darkblue}{HTML}{365F91}
\definecolor{blue}{HTML}{4F81BD}
\definecolor{darkdarkblue}{HTML}{243F60}
\definecolor{codebg}{HTML}{F4F6F8}

\lstset{
  basicstyle=\ttfamily\scriptsize,
  backgroundcolor=\color{codebg},
  frame=single,
  rulecolor=\color{lightgray},
  breaklines=true,
  columns=fullflexible,
  showstringspaces=false,
  captionpos=b,
  xleftmargin=4pt,
  framexleftmargin=4pt,
}

% Section Style
\titleformat{\section}
  {\normalfont\LARGE\upshape\color{darkblue}}{\thesection.}{1em}{}

% Subsection Style
\titleformat{\subsection}
  {\normalfont\Large\upshape\color{blue}}{\thesubsection.}{1em}{}[{\color{black}\titlerule[0.2pt]}]

% Subsubsection Style
\titleformat{\subsubsection}
  {\normalfont\large\upshape\color{blue}}{\thesubsubsection.}{1em}{}[{\color{black}\titlerule[0.2pt]}]

% Paragraph Style
\titleformat{\paragraph}
{\normalfont\normalsize\upshape\color{blue}}{\theparagraph}{1em}{}[{\color{black}\titlerule[0.2pt]}]
\titlespacing*{\paragraph}
{0pt}{3.25ex plus 1ex minus .2ex}{1.5ex plus .2ex}

% Subparagraph Style
\titleformat{\subparagraph}
{\normalfont\normalsize\upshape\itshape\color{darkdarkblue}}{\thesubparagraph}{1em}{}
\titlespacing*{\subparagraph}
{0pt}{3.25ex plus 1ex minus .2ex}{1.5ex plus .2ex}

\pagestyle{myheadings}                          % Own header
\pagestyle{fancy}                               % Own style
\fancyhf{}                                      % Clear header, footer

\setlength{\headheight}{30pt}                   % Set header hight
\renewcommand{\headrulewidth}{0.5pt}            % Top line
\renewcommand{\footrulewidth}{0.5pt}            % Bottom line

\fancyhead[L]{\includegraphics[width=3cm]{guidelines/univienna-logo.eps}} % Header left
\fancyhead[C]{}                                 % Header center
\fancyhead[R]{University of Vienna}             % Header right
\fancyfoot[C]{}                                 % Footer center
\fancyfoot[R]{\thepage/\pageref{LastPage}}      % Footer right


\newcommand{\studentone}{Student 1: Golovanov, Ilja, 12133820}
\newcommand{\studenttwo}{Student 2: Togizbayev, Yernaz, 01429473}

\setcounter{tocdepth}{2}

%%% Begin document %%%
\begin{document}

\begin{figure}                                  % Include Logo
\flushleft
\includegraphics[width=0.4\textwidth]{guidelines/univienna-logo.eps}
\end{figure}

% Set title, author, date
\title{052400-1 VU Information Management and \\Systems Engineering (2025S)\\[1em] { Milestone 2\\[1em]
}}
\author{
\Large{Group 05}\\[1em]
\studentone\ \\[1em]
\studenttwo\ \\[1em]
}

\date{[23.08.2026], Vienna}

\maketitle                                      % Include contents
\newpage
\tableofcontents
\newpage

%%%%%%%%%%%%%%%%%%%%%%%%%%%%%%%%%%%%%%%%%
\section{Milestone 1 - Recap and Revisions}
%%%%%%%%%%%%%%%%%%%%%%%%%%%%%%%%%%%%%%%%%
\subsection{Recap: ER-Diagram }

The conceptual model is unchanged in structure since Milestone 1. Nine entities,
one weak entity, one unary relationship, two m:n relationships and two IS-A
hierarchies. The attributes listed in Section~\ref{sec:revisions} were added
during implementation; no entity or relationship was added or removed.

\begin{figure}[h]
\centering
\includegraphics[width=\textwidth]{docs/ER-diagram/er_diagram.jpg}
\caption{Boat Rental ER diagram (Chen notation)}
\end{figure}

\subsection{Revisions and Changes}
\label{sec:revisions}

No entities or relationships changed. Seven attributes were added, each of them
because the implementation needed a fact the model did not record:

\begin{itemize}
  \item \textbf{Boat.DailyRate} \texttt{DECIMAL(10,2)}. A rental business
        without a price is not one. Nullable, because a manager can add a boat
        before deciding what to charge for it; an unpriced boat is hidden from
        clients rather than offered at nothing.
  \item \textbf{Rental.TotalAmount} \texttt{DECIMAL(10,2)}. What the client was
        charged, frozen when the booking is made. Recomputing it from
        \texttt{DailyRate} would let a later price change rewrite what somebody
        already paid.
  \item \textbf{Rental.StartTime} \texttt{TIME}. A charter starts at a handover
        time, 09:00 by default. Kept separate from \texttt{RentalDate} rather
        than widening it to a datetime, because the composite primary key, the
        rental URLs and the overlap check all key off a date.
  \item \textbf{Rental.CreatedAt} \texttt{DATETIME}. An unpaid booking is a hold
        with a deadline, and the deadline is measured from when the booking was
        made. Without it, a booking made a minute ago is indistinguishable from
        one that lapsed overnight.
  \item \textbf{Rental.CancelledAt} and \textbf{Rental.RefundedAt}
        \texttt{DATETIME}. A paid charter that is called off keeps its row,
        because that row is the only record the client was charged. These two
        say when the refund became owed and when it was returned.
  \item \textbf{Catamaran.HasJacuzzi} \texttt{BOOLEAN}. The MS1 model gave this
        to yachts only, but catamarans have them too, and the booking page
        states the answer either way.
\end{itemize}

The relational structure itself is otherwise as submitted: \texttt{Rental}
remains the weak entity with the composite key
\texttt{(ClientID, BoatID, RentalDate)}, \texttt{Manager.SupervisorID} remains
the unary relationship, and \texttt{Supervises} and \texttt{Maintains} remain
the two m:n tables.

%%%%%%%%%%%%%%%%%%%%%%%%%%%%%%%%%%%%%%%%%
\section{Milestone 2}
%%%%%%%%%%%%%%%%%%%%%%%%%%%%%%%%%%%%%%%%%
\subsection{Infrastructure}

Three containers, defined in \texttt{docker-compose.yml} and started together:

\begin{itemize}
  \item \textbf{db} \textendash{} MariaDB 11.3. The relational database and the
        system of record.
  \item \textbf{mongo} \textendash{} MongoDB 7. The document database the data
        is migrated into.
  \item \textbf{backend} \textendash{} Python 3.11 with Flask, serving both the
        pages and the application logic. It waits for both databases to report
        healthy before starting.
\end{itemize}

All dependencies are installed during the image build from
\texttt{backend/requirements.txt}; nothing has to be installed on the host.

\paragraph{Running it}

\begin{lstlisting}[language=bash]
docker compose up --build
\end{lstlisting}

The application is then at \url{http://localhost:5000}. Both published ports are
bound to \texttt{127.0.0.1} (the app on 5000, MariaDB on 3306, MongoDB on 27017),
so nothing is reachable from outside the machine running it: the Flask
development server would otherwise put a Python traceback in front of anyone on
the same network, and the databases would answer with credentials that are in
the compose file.

To stop and reset the relational database completely:

\begin{lstlisting}[language=bash]
docker compose down -v
\end{lstlisting}

The \texttt{-v} is required to drop the volume, which is what makes the SQL init
scripts run again on the next start.

\paragraph{Signing in} There is no password system. \texttt{/login} lists every
client and one click signs you in; \texttt{/manager/login} picks a manager from
a dropdown. Both roles are needed to see the whole application: the booking
flow is a client's, and the fleet, staff and MongoDB pages are a manager's.

\subsection{RDBMS Implementation}
\subsubsection{DB Setup / Data Import / Base Function of Web System}

The tables are created by the SQL in \texttt{database/}, not by the application.
MariaDB's entrypoint runs \texttt{Group05\_\allowbreak Createtable.sql} on
first start,
followed by \texttt{init.sql}, which sources the per-student insert scripts.
The entrypoint does not recurse into subdirectories, which is why
\texttt{init.sql} exists at the top level purely to \texttt{SOURCE} the files in
\texttt{Student1/} and \texttt{Student2/}.

The application layer mirrors that schema in
\texttt{backend/boat\_rental/\allowbreak models.py} as SQLAlchemy models. The two are
maintained by hand and have to agree: there is no migration tool, so a column
added to one must be added to the other.

\paragraph{Data import} \texttt{backend/boat\_rental/generator.py} fills the
database with randomised data: 200 boats spread over every office, ten clients,
employees split between managers and staff, and roughly 340 rentals. It is
triggered by the \textbf{Generate demo data} button, which is on every page and
does not require signing in, so a fresh checkout can be populated before anyone
has logged in.

Two details are deliberate. The generator replaces all boats, people and
rentals, but it does \emph{not} delete the offices: a manager can open an office
in any city, and clearing them would make the only restocking tool destroy every
city anyone had added. It also runs as a single transaction, so a failure
halfway through cannot leave the database emptied but not refilled.

Boats are generated coherently rather than at random: length is drawn per boat
type, and seats, engine power, weight, cabins and the daily rate are derived
from it, so the fleet does not contain a seven-metre yacht that sleeps five.

\paragraph{Base functionality} The web system covers client registration and
sign-in, harbour and date search, booking with a simulated card payment,
cancellation, and a client logbook. On the manager side it covers full CRUD for
boats, offices and employees, the two m:n assignment relations, cancellation of
any client's charter, and the MongoDB migration and reports described in
Section~2.3.

\section*{Student 1 (Ilja Golovanov)}
\subsubsection{Implementation of the Use Case and the Analytics Report}

\paragraph{Use case} A client books a boat. They choose a harbour and a date
range; the system shows the boats in that harbour, and the client picks one and
pays for the charter.

The whole fleet of that harbour is listed rather than only what is free. Boats
that are already booked or under maintenance are shown greyed out, labelled with
the reason, and carry no radio button at all, so they cannot be selected or
submitted. A page showing three boats reads as a company with three boats;
showing the taken ones tells the client the harbour is worth trying on other
dates.

\paragraph{Analytics report} \emph{Which boats are free in a given harbour over
a given date range, and how much of the fleet that represents.} It is the
\textbf{Availability} page. Alongside the list it shows the fleet size and how
many boats are not under maintenance, so the free count has something to be
measured against.

\paragraph{How the report changes} Booking a boat removes it from the report for
any overlapping date range. The report and the booking page share one overlap
predicate in the code, so what the report calls free and what the booking form
accepts cannot disagree.

\paragraph{Changes since MS1} Two conditions were added to the query as the
implementation exposed them. A boat with no \texttt{DailyRate} is excluded, since
it would otherwise reach the checkout with a total of nothing. A cancelled
rental no longer counts as an overlap, since its row is kept as a payment record
but the boat is free again. The search parameters were narrowed to harbour and
dates; skipper and boat type were not implemented.

\section*{Student 2 (Yernaz Togizbayev)}
\setcounter{subsubsection}{1}
\subsubsection{Implementation of the Use Case and the Analytics Report}

\paragraph{Use case} A manager hires a staff member. The manager fills in the
employee's personal and work details and chooses the role. The application
writes the \texttt{Employee} row and then the matching \texttt{Staff} or
\texttt{Manager} row, and the new staff member is placed under a manager's
supervision through the \texttt{Supervises} relation.

This exercises the IS-A hierarchy directly. Changing an employee's role later
moves them between the \texttt{Staff} and \texttt{Manager} tables, deleting the
old subclass row and writing the new one, rather than trying to update across
them.

\paragraph{Analytics report} \emph{Which staff a given manager supervises, with
each one's shift, duty status and office.} It is the \textbf{Supervision} page.
The SQL joins \texttt{Supervises} to \texttt{Staff}, to \texttt{Employee} twice
(once for the staff member and once for the manager), to \texttt{Manager} and to
\texttt{Office}.

\paragraph{How the report changes} Hiring a staff member and assigning them to a
manager adds a row to that manager's report. Measured on the running system:
before the use case, manager \texttt{E8} supervised three staff; after hiring
Lena Karras and assigning her, the same report returned four.

\paragraph{Changes since MS1} The use case is as designed. One thing the
implementation forced: removing an employee is not a single delete. Three
foreign key families point at them \textendash{}
\texttt{Manager.SupervisorID}, \texttt{Supervises} and \texttt{Maintains}
\textendash{} and all three have to be detached first, including when a role
change deletes the old subclass row. The MS1 design did not account for that.

\pagebreak

\subsection{NoSQL Implementation}

\subsubsection{NoSQL Re-Design}

The relational schema spreads the application over ten tables. Answering either
use case there means joining four or five of them. The document model instead
stores each use case's answer in the shape it is asked for, in three
collections.

\paragraph{offices} One document per harbour, with its boats embedded, and each
boat's rentals embedded inside the boat. Student~1's report becomes a single
aggregation over one collection.

\begin{lstlisting}
{
  "_id": "O5",
  "city": "Dubrovnik",
  "country": "Croatia",
  "street": "Obala Stjepana Radica, 2",
  "zip": "20000",
  "boats": [
    {
      "boatId": "B27",
      "manufacturer": "Lagoon",
      "type": "catamaran",
      "length": 15.2,
      "seats": 10,
      "availabilityStatus": "Available",
      "dailyRate": "1120.00",
      "jacuzzi": true,
      "rentals": [
        { "clientId": "C3", "boatId": "B27",
          "rentalDate": "2026-09-04", "rentalEndDate": "2026-09-09",
          "paymentStatus": "PAID", "totalAmount": "5600.00" }
      ]
    }
  ]
}
\end{lstlisting}

\paragraph{managers} One document per manager, with the staff they supervise
embedded. Student~2's report becomes a single document read.

\begin{lstlisting}
{
  "_id": "E8",
  "employeeId": "E8",
  "name": { "first": "Jamie", "last": "Miller" },
  "email": "jamie.miller@boat-rental.com",
  "office": { "officeId": "O1", "city": "Mykonos", "country": "Greece" },
  "department": "HR",
  "managementLevel": "Senior",
  "supervisorId": null,
  "supervisedStaff": [
    { "employeeId": "E2",
      "name": { "first": "Avery", "last": "Smith" },
      "email": "avery.smith@boat-rental.com",
      "office": { "officeId": "O1", "city": "Mykonos", "country": "Greece" },
      "workShift": "Evening", "isOnDuty": true }
  ]
}
\end{lstlisting}

\paragraph{clients} One document per client with their own rentals embedded,
which is the client logbook.

\paragraph{Keys, and what happened to them} Primary keys are kept, as
\texttt{\_id}: \texttt{OfficeID}, \texttt{ManagerID} and \texttt{ClientID} are
meaningful identifiers that the application already uses in URLs, so replacing
them with generated ObjectIds would add a lookup and gain nothing. Foreign keys
mostly disappear: \texttt{Boat.OfficeID} is implied by the boat being inside its
office document, and \texttt{Supervises} stops being a table at all because the
relationship is the array. The keys that survive are the ones crossing a
document boundary. A rental keeps \texttt{clientId} and \texttt{boatId}, because
it is embedded in two places and each copy needs to name the other side.

The IS-A hierarchies are flattened. \texttt{Employee} is folded into each
manager and staff sub-document, and the boat subtypes become a \texttt{type}
field plus the attributes that subtype has. Nothing in the application asks for
"an employee" or "a boat" without already knowing which kind it is, so the
parent table earns nothing as a separate document.

\paragraph{Changes compared to the individual MS1 designs} Student~1's MS1
design embedded rentals in the client. Student~2's embedded supervised staff in
the manager. Both survive, but neither on its own serves the other's use case,
so the team design keeps both and adds the \texttt{offices} collection with
boats embedded. That is what makes Student~1's report a single-collection query:
his MS1 design would have needed a scan of every client's rentals to find out
whether a boat was free.

The cost is that a rental is stored twice, once under its boat and once under
its client. That is the trade the document model asks for here: both reads
become single-document, and the duplication costs nothing at runtime because
only the migration writes to MongoDB.

\subsubsection{NoSQL DB Setup and Data Migration}

MongoDB 7 runs as the \texttt{mongo} service, with a named volume for its data
and a healthcheck the backend waits on. The application talks to it with
\texttt{pymongo}; all the MongoDB code is in
\texttt{backend/boat\_rental/nosql.py}.

The migration is started from the \textbf{MongoDB} page in the manager
navigation, with the \textbf{Migrate to MongoDB} button. It clears all three
collections before writing, so running it twice produces the same result rather
than duplicating anything.

It is a one-way copy. MariaDB stays the system of record and nothing is written
back to it, and the application never writes to both databases at once: normal
use writes only to MariaDB, and MongoDB is refreshed by re-running the
migration. The data is read through the existing SQLAlchemy models rather than
with fresh SQL, so there is a single definition of what a boat or a rental is.

Rentals are grouped by boat and by client in one pass before the documents are
assembled, so the migration reads each table once rather than querying per
document. On the current demo data one run migrates:

\begin{center}
\begin{tabular}{lr}
\toprule
Collection / embedded & Documents \\
\midrule
offices    & 20 \\
managers   & 4 \\
clients    & 11 \\
boats (embedded in offices)      & 200 \\
rentals (embedded, twice each)   & 323 \\
\bottomrule
\end{tabular}
\end{center}

Semantics are preserved rather than approximated. Money is written as a string,
not a double: the columns are \texttt{DECIMAL(10,2)} and a float cannot hold
cents exactly, so the same reasoning that chose the SQL type chooses the BSON
one. Dates are ISO-8601 strings, which compare correctly as strings and so can
be used directly in the range conditions of the availability query.

\section*{Student 1 (Ilja Golovanov)}
\subsubsection{NoSQL Use Case and Analytics Report}

The use case and the report were reimplemented against MongoDB and are on the
same \textbf{MongoDB} page. Selecting a harbour and a date range runs the
aggregation below and lists the boats it returns.

The results match the relational implementation. This was checked rather than
assumed: for all 20 harbours in the demo data, over the same date range, the SQL
query and the aggregation return exactly the same set of boat identifiers. The
same holds after the use case is triggered: booking a boat in Dubrovnik for a
window that previously had ten free boats leaves nine, and the booked boat is
the one missing from the list.

\subsubsection{NoSQL Analytics Report Statement}

\begin{lstlisting}
db.offices.aggregate([
  { $match: { city: "Dubrovnik" } },
  { $unwind: "$boats" },
  { $match: { "boats.availabilityStatus": "Available",
              "boats.dailyRate": { $ne: null } } },
  { $addFields: {
      clashes: { $filter: {
        input: "$boats.rentals", as: "r",
        cond: { $and: [
          { $lt: [ "$$r.rentalDate", "2026-08-30" ] },
          { $or: [ { $eq: [ "$$r.rentalEndDate", null ] },
                   { $gt: [ "$$r.rentalEndDate", "2026-08-23" ] } ] },
          { $ne: [ "$$r.paymentStatus", "CANCELLED" ] } ] } } } } },
  { $match: { clashes: { $size: 0 } } },
  { $project: { _id: 0, boatId: "$boats.boatId",
                manufacturer: "$boats.manufacturer", type: "$boats.type",
                seats: "$boats.seats", length: "$boats.length",
                jacuzzi: "$boats.jacuzzi", dailyRate: "$boats.dailyRate",
                city: "$city", country: "$country" } },
  { $sort: { manufacturer: 1, boatId: 1 } }
])
\end{lstlisting}

\paragraph{How the design affects it} The whole query stays inside one
collection. The relational version joins \texttt{Boat} to \texttt{Office} and
then tests a correlated \texttt{NOT EXISTS} against \texttt{Rental}; here the
harbour is selected by an indexed equality, and the rentals that would clash are
already inside the boat. The overlap test itself is unchanged in meaning, it has
just moved from a subquery into a \texttt{\$filter} over an embedded array.

\paragraph{Compromises} Two, both accepted knowingly. First, the same rental
document exists under the boat and under the client, so a booking would have to
be written twice. That is tolerable here because nothing but the migration
writes to MongoDB, but it would be the first thing to revisit if the application
started writing to the document database directly. Second, \texttt{\$unwind}
expands one document per boat, so a harbour with a very large fleet does more
work in memory than a row-per-boat table would. The mitigation is that the
\texttt{\$match} on \texttt{city} runs first and is index-backed, so only the
selected harbour is ever unwound. A harbour holding thousands of boats would
push the document toward MongoDB's 16\,MB limit and would be the point to switch
the boats to their own collection with a reference back to the office.

\subsubsection{NoSQL Indexing}

The report selects a harbour by name, so \texttt{offices.city} is indexed:

\begin{lstlisting}
db.offices.createIndex({ city: 1 }, { name: "city_idx" })
\end{lstlisting}

Measured with \texttt{explain("executionStats")} on the demo data, running the
aggregation above with the index dropped and then in place:

\begin{center}
\begin{tabular}{lrrrl}
\toprule
Index & Docs examined & Keys examined & Returned & Plan \\
\midrule
dropped  & 20 & 0 & 1 & COLLSCAN \\
in place & 1  & 1 & 1 & IXSCAN \\
\bottomrule
\end{tabular}
\end{center}

Without the index every harbour document is fetched and discarded; with it,
exactly the one needed is read. The ratio is the number of harbours, so it grows
as the company opens offices. The absolute times are both around a millisecond
at this data size, which is the honest result: the index matters for the number
of documents touched, not for wall-clock time on 20 documents.

\section*{Student 2 (Yernaz Togizbayev)}
\setcounter{subsubsection}{2}
\subsubsection{NoSQL Use Case and Analytics Report}

The use case writes to MariaDB, and the report was reimplemented against
MongoDB on the \textbf{MongoDB} page: selecting a manager shows the staff that
manager supervises, with shift, duty status and office.

The results match the relational implementation, checked across all managers
that supervise anyone: the same staff identifiers come back from both databases.
The report also changes with the use case in the same way. Manager \texttt{E8}
supervised three staff; after hiring Lena Karras, assigning her to \texttt{E8}
and re-running the migration, the MongoDB report returned four, exactly as the
SQL report did.

\subsubsection{NoSQL Analytics Report Statement}

\begin{lstlisting}
db.managers.find({ _id: "E8" })
\end{lstlisting}

\paragraph{How the design affects it} This is the clearest gain of the two.
The relational report joins five tables: \texttt{Supervises},
\texttt{Staff}, \texttt{Employee} (twice), \texttt{Manager} and \texttt{Office}.
Embedding the supervised staff inside the manager reduces it to fetching one
document by its primary key, which MongoDB serves from the \texttt{\_id} index
with no additional index required. Each staff sub-document already carries the
name, shift, duty flag and office, so there is nothing left to look up
afterwards.

\paragraph{Compromises} Embedding means staff details are duplicated wherever a
person appears, so a change of office or shift has to be written in more than
one place. That is acceptable in this application because employee details
change rarely and the report is read often, and because the migration rebuilds
the documents wholesale rather than patching them.

The real cost is directional. The design answers "who does this manager
supervise" in one read, but the reverse question, "who supervises this staff
member", has no natural home: without help it means scanning every manager
document and searching its array. That is what the index in the next section
exists for. A design that instead kept staff in their own collection with a
\texttt{managerId} reference would make the reverse question trivial and the
report a join, which is the wrong trade for this use case, since the report is
what the application actually runs.

\subsubsection{NoSQL Indexing}

The report itself needs no new index, because it is a lookup on \texttt{\_id}.
The reverse lookup does, so the array field is indexed. MongoDB indexes each
element of an embedded array separately, which is what makes this work at all:

\begin{lstlisting}
db.managers.createIndex({ "supervisedStaff.employeeId": 1 },
                        { name: "supervised_staff_idx" })
\end{lstlisting}

Measured with \texttt{explain("executionStats")} on the reverse lookup:

\begin{lstlisting}
db.managers.find({ "supervisedStaff.employeeId": "E2" })
\end{lstlisting}

\begin{center}
\begin{tabular}{lrrrl}
\toprule
Index & Docs examined & Keys examined & Returned & Plan \\
\midrule
dropped  & 4 & 0 & 1 & COLLSCAN \\
in place & 1 & 1 & 1 & IXSCAN + FETCH \\
\bottomrule
\end{tabular}
\end{center}

The collection scan reads every manager document and searches each
\texttt{supervisedStaff} array; the index goes straight to the one that contains
the staff member. With four managers in the demo data the saving is small in
absolute terms, and it would be dishonest to present it as a speed-up at this
size. What the numbers show is the shape of the change: the scan is proportional
to the number of managers, the indexed lookup is not.

Writing this section also caught a real defect. The index was first created on
\texttt{supervisedStaff.staffId}, but the migration writes the field as
\texttt{employeeId}, so the index matched nothing and
\texttt{explain} reported \texttt{nReturned: 0} with the index apparently in
use. The field name was corrected and the figures above are from the corrected
version.

\end{document}
